\section{Scope and Applicability of the Review and Framework}

The review's conclusions concern the mechanisms and operating conditions represented in the cited literature. Differences in platform, threat model, workload, and evaluation outcome limit direct performance comparisons across studies. The synthesis supports architectural comparisons and conditional design recommendations; it does not establish the prevalence of particular methods or provide pooled estimates of their effectiveness.

UAV security, IoT reputation, road-vehicle trust, industrial provenance, and permissioned enterprise ledgers provide candidate mechanisms for cross-domain transfer. Their threat distributions, time scales, connectivity, energy, and governance define the conditions for comparison with air--surface--underwater operations. The synthesis identifies these transfer assumptions at the conceptual level. Target-environment validation is required to establish performance under the acoustic conditions and operational practices of a particular maritime deployment.

OAL, the dependency relation set, CTG loop, evidence profile, hypotheses, and CDUTP form a conceptual architecture for mission-trust analysis. Integrated implementation, formal composition analysis, and mission experiments provide the evaluation pathway for safety, calibration, availability, and audit-performance claims. Evaluation of the four object types should address coverage and consistent application across settings. Practical graph implementations require rules for abstraction, correlation, privacy, evidence expiry, cycles, and conflicting claims, together with measured computation and communication costs. The general assessment $T_o(t,c,m)$ accommodates several representations; a probabilistic implementation requires a defined event and horizon, an explicit statistical model, and calibration under relevant operating conditions.

The evidence categories organize complementary forms of support, with scope and rigor assessed for each claim. Tests of the research hypotheses require concrete comparators, loss functions, and scenario distributions to determine architectural trade-offs and effect sizes. Development of the CDUTP proposal entails a normative schema, interoperable implementations, and a conformance test suite. Formal analysis, controlled studies, field evidence, and multi-party governance exercises provide complementary support for assessing these contributions.

Finally, the public-evidence focus supports reproducible analysis of assurance, resilience, audit, and responsible governance in civil and publicly documented defense-support settings. Sensitive operational details remain within the governance and disclosure controls of the organizations that hold them. Future authorized studies can evaluate the framework under additional contested-environment and platform-specific conditions.
